%% 4-1 きまりから考えを積み上げる：定義・公理・定理（第13回：組 V4-1。ページ制限なし。偶数ページ）
%% 証明は「三角形の内角の和は180°」の1題だけ。二等辺三角形・平行四辺形の性質は対称で説明（証明は省く）。作図の手順は省く。
%% 「定義・公理・定理・法則」の使い方はこの章で決め、ほかの章も同じに使う（速さ＝距離÷時間は定義、内角の和180°は定理、F＝ma は法則）。
%% 合同の記号 ≡ と相似 ∽ は 4-3 の記号カード。この章では「合同」のことばだけを使う。
\makeatletter
%% ---- この章だけで使う部品（頭文字 \bkivb…） ----
% 角のしるし（おうぎ形をぬる）  \bkivbang{(頂点)}{はじめの向き}{終わりの向き}{半径}{色}
\newcommand\bkivbang[5]{%
  \fill[#5!30] #1 -- ($#1+(#2:#4)$) arc (#2:#3:#4) -- cycle;
  \draw[#5!85!black, line width=0.6pt] ($#1+(#2:#4)$) arc (#2:#3:#4);}
% 直角のしるし  \bkivbright{(頂点)}{1本めの辺の向き}{大きさ}
\newcommand\bkivbright[3]{\draw[bkInk, line width=0.5pt] ($#1+(#2:#3)$) -- ++({#2+90}:#3) -- ++({#2+180}:#3);}
% 点と名前  \bkivbpt{(x,y)}{名前}{名前の向き}
\newcommand\bkivbpt[3]{\fill[bkInk] #1 circle (0.8);\node[anchor=#3, font=\sffamily\small, inner sep=0pt, outer sep=1.6mm] at #1 {#2};}
% 式を、まん中に1行で置く  \bkivbfline{…}
\newcommand\bkivbfline[1]{\par\addvspace{3pt}{\centering #1\par}\addvspace{3pt}\noindent\ignorespaces}
% 説明の段の見出し（丸数字＋太字の1行）  \bkivbstep{番号}{見出し}
\newcommand\bkivbstep[2]{\par\addvspace{6pt}\Needspace{4\baselineskip}\noindent
  {\bkcircled[bkOrange]{#1}\enspace\sffamily\bfseries\color{bkOrange!75!black}#2}\par\nobreak\vspace{2pt}}
% 囲みの中の小見出し（それだけで1行）  \bkivbmidashi{見出し}
\newcommand\bkivbmidashi[1]{\par\addvspace{4pt}\Needspace{7\baselineskip}\noindent\textbf{#1}\par\noindent\ignorespaces}
% 等しい長さのしるし（辺のまん中に短い線を1本か2本）  \bkivbtick{(P)}{(Q)}{1|2}
\newcommand\bkivbtick[3]{%
  \coordinate (kivbM) at ($#1!0.5!#2$);
  \ifnum#3=1
    \draw[bkInk, line width=0.6pt] ($(kivbM)!1.6mm!90:#2$) -- ($(kivbM)!1.6mm!-90:#2$);
  \else
    \coordinate (kivbA) at ($(kivbM)!0.6mm!#2$);\coordinate (kivbB) at ($(kivbM)!-0.6mm!#2$);
    \draw[bkInk, line width=0.6pt] ($(kivbA)!1.6mm!90:#2$) -- ($(kivbA)!1.6mm!-90:#2$);
    \draw[bkInk, line width=0.6pt] ($(kivbB)!1.6mm!90:#2$) -- ($(kivbB)!1.6mm!-90:#2$);
  \fi}
\makeatother
%% 証明の囲み
\newtcolorbox{bkivbshoumei}[1][]{enhanced, unbreakable, colframe=bkTheme!70!black, colback=white, boxrule=1pt, arc=2mm,
  left=3mm, right=3mm, top=2mm, bottom=2mm, fonttitle=\sffamily\bfseries\small, coltitle=bkOn, colbacktitle=bkTheme,
  title={#1}}

\bkchapter{4}{1}{きまりから考えを積み上げる：\\定義・公理・定理}{\bklevelrange{2}{3}}{きまりから考えを積み上げる：定義・公理・定理}
\bkmichishirube{1-19,1-20,3-2}{4-6,4-7,6-15}

\begin{goal}
\begin{itemize}[label={\color{\bkTT}●}, leftmargin=1.6\zw, itemsep=0pt, topsep=1pt]
\item \bklevel{2}\enspace 記号 $\angle$・$\triangle$・$\perp$・$\parallel$ を使い、平行線の同位角・錯角から角の大きさを求められる。
\item \bklevel{3}\enspace 三角形の内角の和が180°になるわけを証明でき、多角形の内角の和・外角の和を求められる。
\item \bklevel{3}\enspace 定義・公理・定理・法則のちがいを説明でき、合同条件・逆・反例を使える。
\end{itemize}
\end{goal}

\begin{hatena}
三角形の3つの角を分度器ではかって、たしてみました。どれも、だいたい180°です。
\begin{bkfigure}
\begin{tikzpicture}[x=1mm, y=1mm]
  \foreach \sx/\ax/\ay/\bx/\cx/\cy in {0/0/0/32/12/22, 42/0/0/36/26/18, 88/0/0/26/4/26} {
    \begin{scope}[shift={(\sx,0)}]
      {\filldraw[fill=white, draw=bkBlue!80!black, line width=0.9pt, line join=round] (\ax,\ay) -- (\bx,0) -- (\cx,\cy) -- cycle;}
    \end{scope}}
  \foreach \x/\y/\t in {6.9/4.1/61°, 24.7/3.2/48°, 13/14.1/72°, 53.5/3.6/35°, 71.1/4.1/61°, 66.2/10.2/83°, 94.1/5.2/81°, 106.7/3.4/50°, 94.2/18.3/49°} {
    \node[font=\sffamily\footnotesize, inner sep=0.5pt] at (\x,\y) {\t};}
  \node[anchor=north, font=\sffamily\footnotesize, inner sep=0pt] at (16,-2.5) {$61+48+72=181$};
  \node[anchor=north, font=\sffamily\footnotesize, inner sep=0pt] at (60,-2.5) {$35+61+83=179$};
  \node[anchor=north, font=\sffamily\footnotesize, inner sep=0pt] at (101,-2.5) {$81+50+49=180$};
\end{tikzpicture}
\end{bkfigure}
でも、ぴったり180°ではないこともあります。何回はかっても、世界じゅうの三角形をはかることはできません。それなのに、「三角形の角をたすと、\textbf{必ず}180°」と言えるのは、なぜかな？
\end{hatena}

\noindent はかった値には、いつも少しずれがあります（2-6）。181°や179°になったのは、はかるときのずれでしょうか。それとも、ほんとうに180°でない三角形があるのでしょうか。はかるだけでは、これを決められません。この章では、「はからずに、考えを積み上げて、必ず正しいと言う方法」を学びます。これを\textbf{証明}（しょうめい）といいます。

%% ---------------- やってみよう・思い出そう ----------------
\begin{yattemiyou}[三角形の角を、ちぎって並べよう]
\noindent\begin{minipage}[c]{0.5\linewidth}
\begin{enumerate}[label={\bkcircled[bkGreen]{\arabic*}}, leftmargin=2.4\zw, itemsep=1pt, topsep=0pt]
\item 紙に、好きな形の三角形を1つかいて、はさみで切りとろう。
\item 3つの角に、ちがう色をぬろう。
\item 3つの角を、手でちぎりとろう。
\item ちぎった3つの角の頂点を1つの点に集めて、すきまなく並べよう。どんな形になるかな？
\item ほかの形の三角形でも、ためしてみよう。
\end{enumerate}
\end{minipage}\hfill
\begin{minipage}[c]{0.47\linewidth}\centering
\begin{tikzpicture}[x=1mm, y=1mm]
  % 左：三角形（A=63.4°、B=45°、C=71.6°）
  \bkivbang{(0,0)}{0}{63.43}{6}{bkRed}
  \bkivbang{(30,0)}{135}{180}{6}{bkBlue}
  \bkivbang{(10,20)}{243.43}{315}{6}{bkOrange}
  \draw[bkInk, line width=0.8pt, line join=round] (0,0) -- (30,0) -- (10,20) -- cycle;
  \draw[bk link] (25,14) -- (34,14);
  % 右：一直線に並べる
  \begin{scope}[shift={(48,2)}]
    \bkivbang{(0,0)}{0}{63.43}{7}{bkRed}
    \bkivbang{(0,0)}{63.43}{135}{7}{bkOrange}
    \bkivbang{(0,0)}{135}{180}{7}{bkBlue}
    \draw[bkInk, line width=0.8pt] (-9,0) -- (9,0);
    \node[anchor=north, font=\sffamily\footnotesize, inner sep=1pt] at (0,-1.5) {一直線？};
  \end{scope}
\end{tikzpicture}
\end{minipage}
\end{yattemiyou}

\noindent どんな形の三角形でも、3つの角はすきまなく並んで、ちょうど一直線になったはずです。一直線の角は、半回転で180°です（1-20）。でも、紙はちぎるとき少し欠けるし、並べるときにすきまができることもあります。これも、「必ず」とまでは言えません。

\begin{omoidasou}
\bkrecall{1-20}{1回転は360°。半回転（一直線）は180°、直角は90°。}
\bkrecall{1-19}{三角形の3つの角のかどを頂点、まわりの線を辺という。}
\bkrecall{2-6}{はかった値には、いつも少しずれがある。}
\bkrecall{3-2}{点 $(x,\ y)$ は、原点から右へ $x$、上へ $y$ の所。}
\end{omoidasou}

%% ---------------- 数学の道具 ----------------
\bksection{doguu}{数学の道具}

\bksubsection{1}{点と線の名前と、記号}

\label{bk:4-1:kigou}考えを積み上げるには、まず、ことばの意味を全員でそろえる必要があります。人によってことばの意味がちがうと、話がかみ合わないからです。図形でいちばん土台になる「点」と「線」から始めます。

\textbf{点}は、場所だけを表し、大きさはありません。紙にかくときは小さな丸でかきますが、ほんとうは大きさ0です。点には、A、B、Pのような大文字のアルファベットで名前をつけます。

\begin{bkfigure}
\begin{tikzpicture}[x=1mm, y=1mm]
  % 直線
  \draw[bkBlue, line width=1pt] (10,24) -- (50,24);
  \draw[bkBlue, line width=1pt, dash pattern=on 1pt off 1.5pt] (2,24) -- (10,24) (50,24) -- (58,24);
  \bkivbpt{(20,24)}{A}{south}\bkivbpt{(40,24)}{B}{south}
  \node[anchor=west, font=\sffamily\bfseries\small, inner sep=0pt] at (64,24) {直線 AB};
  \node[anchor=west, font=\small, inner sep=0pt] at (84,24) {両方のはしが、どこまでものびる};
  % 線分
  \draw[bkBlue, line width=1pt] (20,12) -- (40,12);
  \bkivbpt{(20,12)}{A}{south}\bkivbpt{(40,12)}{B}{south}
  \node[anchor=west, font=\sffamily\bfseries\small, inner sep=0pt] at (64,12) {線分 AB};
  \node[anchor=west, font=\small, inner sep=0pt] at (84,12) {AとBで終わる。長さがある};
  % 半直線
  \draw[bkBlue, line width=1pt] (20,0) -- (50,0);
  \draw[bkBlue, line width=1pt, dash pattern=on 1pt off 1.5pt] (50,0) -- (58,0);
  \bkivbpt{(20,0)}{A}{south}\bkivbpt{(40,0)}{B}{south}
  \node[anchor=west, font=\sffamily\bfseries\small, inner sep=0pt] at (64,0) {半直線 AB};
  \node[anchor=west, font=\small, inner sep=0pt] at (84,0) {Aから始まり、Bのほうへのびる};
\end{tikzpicture}
\bkfigcaption{直線・線分・半直線。点線は「この先もずっと続く」というしるし。}
\end{bkfigure}

2つの点AとBを通る、まっすぐな線を考えます。両方へどこまでものびる線を\textbf{直線}AB、AからBまでで終わる線を\textbf{線分}（せんぶん）AB、Aから始まってBのほうへどこまでものびる線を\textbf{半直線}（はんちょくせん）ABといいます。紙にかけるのは直線の一部分だけですが、頭の中では、どこまでものびていると考えます。線分ABの長さも、そのまま AB と書きます。たとえば「AB＝\qty{3}{cm}」です。

\bkivbmidashi{垂直と平行}
2本の直線の関係には、とくに大事なものが2つあります。方眼のノートの、たての線と横の線を思いうかべてください。

\begin{bkfigure}
\begin{tikzpicture}[x=1mm, y=1mm]
  % 垂直
  \begin{scope}[shift={(22,12)}, rotate=20]
    \draw[bkBlue, line width=1pt] (-20,0) -- (20,0);
    \draw[bkOrange, line width=1pt] (0,-13) -- (0,13);
    \bkivbright{(0,0)}{0}{2.6}
    \bkivbpt{(-14,0)}{A}{north}\bkivbpt{(14,0)}{B}{north}
    \bkivbpt{(0,-9)}{C}{west}\bkivbpt{(0,9)}{D}{west}
  \end{scope}
  \node[font=\sffamily\bfseries\small] at (22,-8) {$\mathrm{AB}\perp\mathrm{CD}$};
  % 平行
  \begin{scope}[shift={(84,12)}, rotate=12]
    \draw[bkBlue, line width=1pt] (-24,5) -- (24,5);
    \draw[bkOrange, line width=1pt] (-24,-5) -- (24,-5);
    \bkivbpt{(-14,5)}{A}{south}\bkivbpt{(12,5)}{B}{south}
    \bkivbpt{(-12,-5)}{C}{north}\bkivbpt{(14,-5)}{D}{north}
    \draw[bkGray, line width=0.5pt, <->] (-2,-5) -- (-2,5);
    \draw[bkGray, line width=0.5pt, <->] (20,-5) -- (20,5);
  \end{scope}
  \node[font=\sffamily\bfseries\small] at (84,-8) {$\mathrm{AB}\parallel\mathrm{CD}$};
\end{tikzpicture}
\bkfigcaption{左：垂直。直角に交わる。右：平行。どこまでのばしても交わらない（あいだのはばは、どこでも同じ）。}
\end{bkfigure}

\begin{itemize}[leftmargin=1.2\zw, itemsep=2pt, topsep=2pt, label={\color{bkTheme!80!black}●}]
\item 2本の直線が直角に交わるとき、2本は\textbf{垂直}（すいちょく）であるといい、$\mathrm{AB}\perp\mathrm{CD}$ と書きます。方眼のたての線と横の線は垂直です。
\item 同じ平面の上の2本の直線が、どこまでのばしても交わらないとき、2本は\textbf{平行}（へいこう）であるといい、$\mathrm{AB}\parallel\mathrm{CD}$ と書きます。方眼の横の線どうしは平行です。電車のレールも平行です。
\end{itemize}

\bkivbmidashi{角と三角形の書き方}
\par\addvspace{4pt}\noindent\begin{minipage}[c]{0.5\linewidth}
1つの点から出た2本の半直線は、\textbf{角}をつくります。点Bから、半直線BAと半直線BCが出ているとき、この角を $\angle\mathrm{ABC}$ と書き、「角ABC」と読みます。

角の頂点（ちょうてん）Bを、まん中に書くのがきまりです。ほかに角がなくてまぎれないときは、短く $\angle\mathrm{B}$ とも書きます。角の大きさも同じ記号で、$\angle\mathrm{ABC}=40°$ のように書きます。

3つの点A、B、Cを線分で結んだ三角形を、$\triangle\mathrm{ABC}$ と書き、「三角形ABC」と読みます。
\end{minipage}\hfill
\begin{minipage}[c]{0.47\linewidth}\centering
\begin{tikzpicture}[x=1mm, y=1mm]
  \bkivbang{(0,0)}{0}{40}{7}{bkOrange}
  \draw[bkInk, line width=0.9pt] (0,0) -- (23,0);
  \draw[bkInk, line width=0.9pt] (0,0) -- (40:23);
  \bkivbpt{(40:18)}{A}{south east}\bkivbpt{(0,0)}{B}{east}\bkivbpt{(18,0)}{C}{north}
  \node[font=\sffamily\footnotesize, inner sep=0.5pt] at (20:11.5) {40°};
  \node[anchor=north, font=\sffamily\footnotesize, inner sep=0pt] at (11,-6) {$\angle\mathrm{ABC}$（$\angle\mathrm{B}$）};
  \begin{scope}[shift={(32,0)}]
    \draw[bkBlue, line width=0.9pt, line join=round, fill=bkBlue!8] (0,0) -- (20,0) -- (7,16) -- cycle;
    \node[anchor=north east, font=\sffamily\small, inner sep=1.5pt] at (0,0) {B};
    \node[anchor=north west, font=\sffamily\small, inner sep=1.5pt] at (20,0) {C};
    \node[anchor=south, font=\sffamily\small, inner sep=1.5pt] at (7,16) {A};
    \node[anchor=north, font=\sffamily\footnotesize, inner sep=0pt] at (10,-6) {$\triangle\mathrm{ABC}$};
  \end{scope}
\end{tikzpicture}
\end{minipage}

\par\vspace{4pt}
$\triangle$ は、2-3で学んだ $\Delta$（デルタ、変化量）とは別の記号です。$\triangle$ は「三角形」の絵、$\Delta$ はギリシャ文字です。$\triangle\mathrm{ABC}$ は「Δ×ABC」でも「ABCの変化量」でもありません。

\bkkigocards{\bkkigocard{$\angle$・$\triangle$}{角・三角形}{$\angle\mathrm{ABC}$ は頂点Bの角。$\triangle\mathrm{ABC}$ は三角形ABC。$\triangle$ は $\Delta$（変化量）とはちがう}{angle / triangle}{4-1}}
{\bkkigocard{$\perp$・$\parallel$}{垂直・平行}{$\mathrm{AB}\perp\mathrm{CD}$ は直角に交わる。$\mathrm{AB}\parallel\mathrm{CD}$ はどこまでも交わらない}{perpendicular /\linebreak[0] \mbox{parallel}}{4-1}}

\bksubsection{2}{ずらす・回す・裏返す：合同と対称}

\label{bk:4-1:idou}うすい紙を三角形の上に重ねて、形を写しとります。その紙を動かして、ほかの三角形にぴったり重なるか調べてみましょう。紙の動かし方は、3つしかありません。

\begin{bkfigure}
\begin{tikzpicture}[x=1mm, y=1mm]
  % もとの三角形 (0,0)-(14,0)-(4,10) を、灰色の点線でかく
  \foreach \sx in {0,46,92} {\begin{scope}[shift={(\sx,0)}]
    {\draw[bkGray, line width=0.8pt, dashed, line join=round] (0,0) -- (14,0) -- (4,10) -- cycle;}
  \end{scope}}
  % ずらす：右へ16、上へ8
  \draw[bkBlue, line width=1pt, fill=bkBlue!12, line join=round] (16,8) -- (30,8) -- (20,18) -- cycle;
  \draw[bk link, bkBlue] (6,4) -- (22,12);
  \node[anchor=north, font=\sffamily\bfseries\small, inner sep=0pt] at (17,-4) {ずらす};
  \node[anchor=north, font=\sffamily\footnotesize, text=bkGray!60!black, inner sep=0pt] at (17,-9) {（平行移動）};
  % 回す：点Oのまわりに 90°
  \begin{scope}[shift={(46,0)}]
    \fill[bkInk] (18,2) circle (0.9);
    \node[anchor=north, font=\sffamily\small, inner sep=1.5pt] at (18,2) {O};
    \draw[bkOrange, line width=1pt, fill=bkOrange!15, line join=round] (34,20) -- (34,6) -- (26,16) -- cycle;
    \draw[-{Stealth[length=4pt]}, bkOrange!80!black, line width=0.8pt] (18,2) ++(160:9) arc (160:70:9);
    \node[anchor=north, font=\sffamily\bfseries\small, inner sep=0pt] at (17,-4) {回す};
    \node[anchor=north, font=\sffamily\footnotesize, text=bkGray!60!black, inner sep=0pt] at (17,-9) {（回転移動）};
  \end{scope}
  % 裏返す：直線 x=17 で折り返す
  \begin{scope}[shift={(92,0)}]
    \draw[bkGreen!60!black, line width=0.8pt, dash pattern=on 3pt off 1.5pt] (17,-2) -- (17,16);
    \draw[bkGreen!60!black, line width=1pt, fill=bkGreen!15, line join=round] (34,0) -- (20,0) -- (30,10) -- cycle;
    \node[anchor=north, font=\sffamily\bfseries\small, inner sep=0pt] at (17,-4) {裏返す};
    \node[anchor=north, font=\sffamily\footnotesize, text=bkGray!60!black, inner sep=0pt] at (17,-9) {（対称移動）};
  \end{scope}
\end{tikzpicture}
\bkfigcaption{紙を動かす3つの方法。点線がもとの三角形。裏返すときは、緑の点線で折る。}
\end{bkfigure}

どの動かし方でも、形と大きさは変わりません。ずらす・回す・裏返すを組み合わせて、ぴったり重なる2つの図形を、\textbf{合同}（ごうどう）であるといいます。\bkidx{ごうどう}{合同}
\begin{itemize}[leftmargin=1.2\zw, itemsep=2pt, topsep=2pt, label={\color{bkTheme!80!black}●}]
\item 重なり合う頂点・辺・角を、それぞれ「対応する」頂点・辺・角といいます。
\item 合同な図形では、対応する辺の長さは等しく、対応する角の大きさも等しい。
\item $\triangle\mathrm{ABC}$ と $\triangle\mathrm{DEF}$ が合同で、AとD、BとE、CとFが重なるときは、対応する頂点を同じ順に書いて「$\triangle\mathrm{ABC}$ と $\triangle\mathrm{DEF}$ は合同」といいます（合同の記号は4-3で）。
\end{itemize}

\bkivbmidashi{線対称と点対称}
1つの図形が、自分自身にぴったり重なることもあります。

\begin{bkfigure}
\begin{tikzpicture}[x=1mm, y=1mm]
  % 線対称：二等辺三角形
  \bkivbang{(0,0)}{0}{56.3}{6}{bkBlue}
  \bkivbang{(32,0)}{123.7}{180}{6}{bkBlue}
  \draw[bkInk, line width=1pt, line join=round] (0,0) -- (32,0) -- (16,24) -- cycle;
  \draw[bkGreen!60!black, line width=0.9pt, dash pattern=on 3pt off 1.5pt] (16,28) -- (16,-4);
  \bkivbright{(16,0)}{0}{2.4}
  \bkivbtick{(0,0)}{(16,24)}{1}\bkivbtick{(32,0)}{(16,24)}{1}
  \bkivbtick{(0,0)}{(16,0)}{2}\bkivbtick{(16,0)}{(32,0)}{2}
  \node[anchor=south, font=\sffamily\small, inner sep=1.5pt] at (16,28) {対称の軸};
  \node[anchor=north east, font=\sffamily\small, inner sep=1.5pt] at (0,0) {B};
  \node[anchor=north west, font=\sffamily\small, inner sep=1.5pt] at (32,0) {C};
  \node[anchor=south east, font=\sffamily\small, inner sep=1.5pt] at (15.5,24) {A};
  \node[anchor=north, font=\sffamily\bfseries\small, inner sep=0pt] at (16,-7) {線対称（二等辺三角形）};
  % 点対称：平行四辺形
  \begin{scope}[shift={(68,0)}]
    \bkivbang{(0,0)}{0}{57.53}{5}{bkOrange}
    \bkivbang{(46,22)}{180}{237.53}{5}{bkOrange}
    \draw[bkInk, line width=1pt, line join=round] (0,0) -- (32,0) -- (46,22) -- (14,22) -- cycle;
    \draw[bkGray, line width=0.5pt] (0,0) -- (46,22) (32,0) -- (14,22);
    \fill[bkRed] (23,11) circle (1);
    \node[anchor=north, font=\sffamily\small, inner sep=2pt] at (23,10) {O};
    \draw[-{Stealth[length=4pt]}, bkRed, line width=0.8pt] (23,11) ++(110:9) arc (110:270:9);
    \bkivbtick{(0,0)}{(32,0)}{1}\bkivbtick{(14,22)}{(46,22)}{1}
    \bkivbtick{(32,0)}{(46,22)}{2}\bkivbtick{(0,0)}{(14,22)}{2}
    \node[anchor=north, font=\sffamily\bfseries\small, inner sep=0pt] at (23,-7) {点対称（平行四辺形）};
  \end{scope}
\end{tikzpicture}
\bkfigcaption{同じしるしの辺は長さが等しく、同じ色の角は大きさが等しい。}
\end{bkfigure}

\begin{itemize}[leftmargin=1.2\zw, itemsep=2pt, topsep=2pt, label={\color{bkTheme!80!black}●}]
\item 1本の直線で折ると、ぴったり重なる図形を\textbf{線対称}（せんたいしょう）な図形といい、折り目の直線を\textbf{対称の軸}といいます。
\item 1つの点のまわりに180°回すと、ぴったり重なる図形を\textbf{点対称}（てんたいしょう）な図形といい、その点を\textbf{対称の中心}といいます。
\end{itemize}

\textbf{二等辺三角形}（にとうへんさんかくけい）は、「2つの辺の長さが等しい三角形」です。等しい2つの辺がはさむ角を\textbf{頂角}（ちょうかく）、残りの2つの角を\textbf{底角}（ていかく）といいます。二等辺三角形を切りとって、頂点Aと、底辺BCのまん中を通る線で折ると、ABとACがぴったり重なります。だから二等辺三角形は線対称で、次のことがわかります。
\begin{enumerate}[label={\bkcircled[bkDeep]{\arabic*}}, leftmargin=2.4\zw, itemsep=1pt, topsep=2pt]
\item 2つの底角は等しい（$\angle\mathrm{B}=\angle\mathrm{C}$）。折ると重なるから。
\item 対称の軸は、底辺BCを2等分し、BCに垂直に交わる。
\end{enumerate}

\textbf{平行四辺形}（へいこうしへんけい）は、「2組の向かい合う辺が、それぞれ平行な四角形」です。平行四辺形を写しとった紙を、2本の対角線（たいかくせん：となり合わない頂点を結ぶ線分）の交わる点Oにピンでとめて180°回すと、もとの平行四辺形にぴったり重なります。だから平行四辺形は点対称で、次のことがわかります。
\begin{enumerate}[label={\bkcircled[bkDeep]{\arabic*}}, leftmargin=2.4\zw, itemsep=1pt, topsep=2pt]
\item 向かい合う辺の長さは等しい。向かい合う角の大きさも等しい。
\item 2本の対角線は、それぞれのまん中の点で交わる。
\end{enumerate}
（どちらも、紙を折ったり回したりして確かめたことです。正しいことを証明する方法もありますが、この本では省きます。）

\par\addvspace{4pt}\noindent\begin{minipage}[c]{0.4\linewidth}\centering
\begin{tikzpicture}[x=1mm, y=1mm]
  \draw[bkLine, line width=0.3pt, step=5] (-20,-12.5) grid (20,12.5);
  \draw[bkInk, line width=0.6pt] (-21,0) -- (21,0) (0,-13) -- (0,13);
  \draw[bkRed, line width=0.8pt, dash pattern=on 2pt off 1.5pt] (15,5) -- (-15,-5);
  \fill[bkBlue] (15,5) circle (1.1);
  \fill[bkOrange] (-15,-5) circle (1.1);
  \node[anchor=south west, font=\sffamily\footnotesize, fill=white, inner sep=0.8pt] at (15.5,5.5) {$(3,\ 1)$};
  \node[anchor=north east, font=\sffamily\footnotesize, fill=white, inner sep=0.8pt] at (-15.5,-5.5) {$(-3,\ -1)$};
  \node[anchor=north west, font=\sffamily\footnotesize, inner sep=1pt] at (0.3,-0.3) {O};
\end{tikzpicture}
\end{minipage}\hfill
\begin{minipage}[c]{0.57\linewidth}
\textbf{座標で見る点対称}\par
原点Oを中心に、点 $(3,\ 1)$ を180°回します。Oと点を結ぶ線分も、いっしょに半回転します。すると点は、Oをはさんで反対側の、同じ距離の所に来ます。右へ3が左へ3に、上へ1が下へ1になるので、$(-3,\ -1)$ です（3-2）。原点のまわりに180°回すと、\textbf{点 $(x,\ y)$ は $(-x,\ -y)$ に移ります}。180°以外の角だけ回したときの式は、4-8で作ります。
\end{minipage}

\bksubsection{3}{平行線と角：対頂角・同位角・錯角}

\label{bk:4-1:heikou}\bkivbmidashi{対頂角}
\par\addvspace{4pt}\noindent\begin{minipage}[c]{0.4\linewidth}\centering
\begin{tikzpicture}[x=1mm, y=1mm]
  \bkivbang{(0,0)}{125}{195}{7}{bkRed}
  \bkivbang{(0,0)}{15}{125}{6}{bkGray}
  \bkivbang{(0,0)}{-55}{15}{7}{bkRed}
  \draw[bkInk, line width=0.9pt] (195:24) -- (15:24);
  \draw[bkInk, line width=0.9pt] (305:18) -- (125:18);
  \node[font=\sffamily\small, inner sep=0.5pt] at (160:10.5) {$a$};
  \node[font=\sffamily\small, inner sep=0.5pt] at (70:9.5) {$b$};
  \node[font=\sffamily\small, inner sep=0.5pt] at (-20:10.5) {$c$};
\end{tikzpicture}
\end{minipage}\hfill
\begin{minipage}[c]{0.57\linewidth}
2本の直線が交わると、4つの角ができます。向かい合う2つの角（$a$ と $c$）を、\textbf{対頂角}（たいちょうかく）といいます。

分度器ではかると、$a$ は70°、$b$ は110°、$c$ は70°でした。$a$ と $c$ は等しくなっています。
\end{minipage}

\par\vspace{4pt}
これは、はからなくてもわかります。$a$ と $b$ は、合わせると一直線（180°）です。$b$ と $c$ も、合わせると一直線です。
\bkivbfline{$a+b=180°$\qquad$b+c=180°$}
どちらも「$b$ にたすと180°になる角」なので、$a$ と $c$ は同じ大きさです。$b$ が110°なら、$a$ も $c$ も $180-110=70$ で70°です。$b$ が何度でも、同じ理由で $a=c$ になります。\textbf{対頂角は等しい}。

\bkivbmidashi{同位角と錯角}
2本の直線 $\ell$、$m$ に、もう1本の直線 $n$ が交わっています（$\ell$ はエル、$m$ はエム。直線の名前には小文字もよく使う）。交わる点が2つあり、それぞれに4つずつ角ができます。

\begin{bkfigure}
\begin{tikzpicture}[x=1mm, y=1mm]
  % 左：平行でない。右：平行。n は 60° の向き
  \foreach \sx/\mang/\tit in {0/-8/{$\ell$ と $m$ が平行でない}, 70/0/{$\ell\parallel m$}} {
    \begin{scope}[shift={(\sx,0)}]
      {\coordinate (P) at (17.3,30);
      \coordinate (Q) at (5.77,10);
      \bkivbang{(P)}{0}{60}{6}{bkBlue}
      \bkivbang{(P)}{180}{240}{6}{bkOrange}
      \bkivbang{(Q)}{\mang}{60}{6}{bkBlue}
      \draw[bkInk, line width=0.9pt] (-4,30) -- (50,30);
      \draw[bkInk, line width=0.9pt] (Q) ++({\mang+180}:12) -- ++(\mang:56);
      \draw[bkInk, line width=0.9pt] (0,0) -- (23,39.8);
      \node[anchor=west, font=\sffamily\small, inner sep=1pt] at (51,30) {$\ell$};
      \node[anchor=west, font=\sffamily\small, inner sep=1pt] at ($(Q)+(\mang:44.5)$) {$m$};
      \node[anchor=south, font=\sffamily\small, inner sep=1pt] at (23,39.8) {$n$};
      \node[font=\sffamily\small, inner sep=0.5pt] at ($(P)+(30:9.5)$) {$p$};
      \node[font=\sffamily\small, inner sep=0.5pt] at ($(P)+(210:9.5)$) {$r$};
      \node[font=\sffamily\small, inner sep=0.5pt] at ($(Q)+({(\mang+60)/2}:9.5)$) {$q$};
      \node[anchor=north, font=\sffamily\bfseries\small, inner sep=0pt] at (24,-3) {\tit};}
    \end{scope}}
  \node[anchor=north, font=\sffamily\footnotesize, inner sep=0pt] at (24,-9) {$p=60°$、$q=68°$};
  \node[anchor=north, font=\sffamily\footnotesize, inner sep=0pt] at (94,-9) {$p=60°$、$q=60°$};
\end{tikzpicture}
\bkfigcaption{$p$ と $q$ は同位角、$r$ と $q$ は錯角。$p$ と $r$ は対頂角。}
\end{bkfigure}

\begin{itemize}[leftmargin=1.2\zw, itemsep=2pt, topsep=2pt, label={\color{bkTheme!80!black}●}]
\item $p$ と $q$ のように、2つの交わる点で「同じ位置」（どちらも右上）にある角を、\textbf{同位角}（どういかく）といいます。
\item $r$ と $q$ のように、$\ell$ と $m$ のあいだで、$n$ をはさんで「ななめに向かい合う」角を、\textbf{錯角}（さっかく）といいます。\bkidx{さっかく}{錯角}
\end{itemize}

図の左のように、$\ell$ と $m$ が平行でないときは、同位角 $p$ と $q$ の大きさはちがいます。右のように $\ell\parallel m$ のときは、$p$ と $q$ は同じ大きさです。$\ell$ を $n$ にそってまっすぐ下へずらすと、平行な $m$ にぴったり重なり、交わる所の形もまるごと重なるからです。

\begin{taisetsu}
平行な2本の直線に、1本の直線が交わるとき、\par
\quad\textbf{同位角は等しい}。\qquad\textbf{錯角は等しい}。
\end{taisetsu}

錯角が等しいわけは、同位角からわかります。図の右で、$r$ と $p$ は対頂角なので $r=p$。$p$ と $q$ は同位角なので $p=q$。だから $r=q$ です。

\begin{seikaku}
\textbf{反対も言える}\par
反対に、2本の直線に1本の直線が交わってできる同位角（または錯角）が等しいなら、その2本の直線は平行です。三角定規を定規にそってすべらせて平行な線を引けるのは、すべらせても三角定規の角（同位角）が変わらないからです。このように、「〜ならば…」の前と後ろを入れかえたものを「逆」といいます（⑦でくわしく）。
\end{seikaku}

\bksubsection{4}{三角形の内角の和は180°}

\label{bk:4-1:naikaku}三角形の内側にある3つの角を\textbf{内角}（ないかく）といいます。\bkhatenaref で、内角の和をはかると180°に近くなりました。やってみようでは、3つの角が一直線に並びました。この「一直線」を、紙をちぎらずに、頭の中で作ります。

\textbf{アイデア}：頂点Aを通って、辺BCに平行な直線DEを引きます。すると、③の「平行線の錯角は等しい」が使えて、$\angle\mathrm{B}$ と $\angle\mathrm{C}$ を、頂点Aのまわりに「移す」ことができます。

\begin{bkfigure}
\begin{tikzpicture}[x=1mm, y=1mm]
  \bkivbang{(0,0)}{0}{54.46}{7}{bkBlue}
  \bkivbang{(66,0)}{141.34}{180}{8}{bkOrange}
  \bkivbang{(24,33.6)}{180}{234.46}{7}{bkBlue}
  \bkivbang{(24,33.6)}{234.46}{321.34}{5}{bkRed}
  \bkivbang{(24,33.6)}{321.34}{360}{8}{bkOrange}
  \draw[bkInk, line width=1pt, line join=round] (0,0) -- (66,0) -- (24,33.6) -- cycle;
  \draw[bkGreen!60!black, line width=1pt] (-8,33.6) -- (72,33.6);
  \draw[-{Stealth[length=4pt]}, bkBlue!80!black, line width=0.8pt] (7,9) to[bend left=35] (13,29);
  \draw[-{Stealth[length=4pt]}, bkOrange!80!black, line width=0.8pt] (57,5) to[bend right=35] (36,30.5);
  \node[anchor=north east, font=\sffamily\small, inner sep=1.5pt] at (0,0) {B};
  \node[anchor=north west, font=\sffamily\small, inner sep=1.5pt] at (66,0) {C};
  \node[anchor=south, font=\sffamily\small, inner sep=1.5pt] at (24,34) {A};
  \node[anchor=south, font=\sffamily\small, inner sep=1.5pt] at (-6,34) {D};
  \node[anchor=south, font=\sffamily\small, inner sep=1.5pt] at (70,34) {E};
  \node[anchor=west, text width=46mm, align=left, font=\small, inner sep=0pt] at (80,18)
    {緑の線DEは、辺BCに平行。\\[3pt]
     青の角（$\angle\mathrm{B}$）と、オレンジの角（$\angle\mathrm{C}$）が、Aのまわりに移る。\\[3pt]
     赤の角（$\angle\mathrm{A}$）とならべると、一直線のDEになる。};
\end{tikzpicture}
\bkfigcaption{アイデア：平行線を1本引いて、3つの角をAのまわりに集める（やってみようの「ちぎって並べる」を頭の中でする）}
\end{bkfigure}

\begin{seikaku}
\textbf{証明の読み方}\par
証明は、「わかっていること（仮定・かてい）」から出発して、1行に1つずつ、理由のはっきりしたことだけを使って、「示したいこと（結論・けつろん）」までたどりつく文章です。どの行にも「なぜそう言えるか」が書いてあるので、読む人は1行ずつ確かめられます。証明は、この本では次の1つだけです。ほかの性質は、図で確かめるだけにします。
\end{seikaku}

\begin{bkivbshoumei}[証明：三角形の内角の和は180°]
\noindent\begin{minipage}[c]{0.44\linewidth}\centering
\begin{tikzpicture}[x=1mm, y=1mm]
  \bkivbang{(0,0)}{0}{59.04}{6}{bkBlue}
  \bkivbang{(50,0)}{136.85}{180}{7}{bkOrange}
  \bkivbang{(18,30)}{180}{239.04}{6}{bkBlue}
  \bkivbang{(18,30)}{239.04}{316.85}{5}{bkRed}
  \bkivbang{(18,30)}{316.85}{360}{7}{bkOrange}
  \draw[bkInk, line width=1pt, line join=round] (0,0) -- (50,0) -- (18,30) -- cycle;
  \draw[bkGreen!60!black, line width=1pt] (-4,30) -- (54,30);
  \node[anchor=north east, font=\sffamily\small, inner sep=1.5pt] at (0,0) {B};
  \node[anchor=north west, font=\sffamily\small, inner sep=1.5pt] at (50,0) {C};
  \node[anchor=south, font=\sffamily\small, inner sep=1.5pt] at (18,30.5) {A};
  \node[anchor=south, font=\sffamily\small, inner sep=1.5pt] at (-2,30.5) {D};
  \node[anchor=south, font=\sffamily\small, inner sep=1.5pt] at (52,30.5) {E};
\end{tikzpicture}
\end{minipage}\hfill
\begin{minipage}[c]{0.53\linewidth}
\textbf{わかっていること（仮定）}\par
$\triangle\mathrm{ABC}$ は三角形。\par\vspace{3pt}
\textbf{示したいこと（結論）}\par
$\angle\mathrm{A}+\angle\mathrm{B}+\angle\mathrm{C}=180°$\par\vspace{3pt}
\textbf{準備}\par
頂点Aを通り、辺BCに平行な直線DEを引く（緑の線）。
\end{minipage}
\par\vspace{4pt}
\begin{enumerate}[label={\bkcircled[bkDeep]{\arabic*}}, leftmargin=2.4\zw, itemsep=2pt, topsep=0pt]
\item DE $\parallel$ BC で、平行線の錯角は等しいので、$\angle\mathrm{DAB}=\angle\mathrm{B}$（\textcolor{bkBlue!80!black}{青}）
\item 同じように、平行線の錯角は等しいので、$\angle\mathrm{EAC}=\angle\mathrm{C}$（\textcolor{bkOrange!80!black}{オレンジ}）
\item D・A・Eは一直線の上にあるので、$\angle\mathrm{DAB}+\angle\mathrm{BAC}+\angle\mathrm{EAC}=180°$
\item ③の式の $\angle\mathrm{DAB}$ を $\angle\mathrm{B}$ に、$\angle\mathrm{EAC}$ を $\angle\mathrm{C}$ におきかえる（①・②）と、\newline $\angle\mathrm{B}+\angle\mathrm{A}+\angle\mathrm{C}=180°$
\end{enumerate}
\hfill だから、三角形の内角の和は180°である。
\end{bkivbshoumei}

証明の中では、ものさしも分度器も使っていません。使ったのは、「平行線の錯角は等しい」と「一直線は180°」の2つだけです。この2つは、とがった三角形でも、平たい三角形でも、どんな三角形でも成り立ちます。だから、結論も、どんな三角形でも成り立ちます。

\bkhatenaref の答え：三角形の角をたすと「必ず」180°と言えるのは、どんな三角形でも成り立つことだけを使って、180°になることを証明できるからです。181°や179°になったのは、はかるときのずれです。


\bkivbmidashi{外角}
\par\addvspace{4pt}\noindent\begin{minipage}[c]{0.46\linewidth}\centering
\begin{tikzpicture}[x=1mm, y=1mm]
  \bkivbang{(0,0)}{0}{50}{6}{bkBlue}
  \bkivbang{(17.77,21.18)}{230}{300}{5}{bkOrange}
  \bkivbang{(30,0)}{120}{180}{4.5}{bkGray}
  \bkivbang{(30,0)}{0}{120}{7}{bkGreen}
  \draw[bkInk, line width=1pt, line join=round] (0,0) -- (30,0) -- (17.77,21.18) -- cycle;
  \draw[bkInk, line width=1pt] (30,0) -- (50,0);
  \node[anchor=north east, font=\sffamily\small, inner sep=1.5pt] at (0,0) {B};
  \node[anchor=north, font=\sffamily\small, inner sep=1.5pt] at (30,-0.5) {C};
  \node[anchor=north, font=\sffamily\small, inner sep=1.5pt] at (50,-0.5) {D};
  \node[anchor=south, font=\sffamily\small, inner sep=1.5pt] at (17.77,21.7) {A};
  \node[font=\sffamily\footnotesize, inner sep=0.5pt] at (25:10) {50°};
  \node[font=\sffamily\footnotesize, inner sep=0.5pt] at ($(17.77,21.18)+(265:9.5)$) {70°};
  \node[font=\sffamily\footnotesize, inner sep=0.5pt] at ($(30,0)+(45:11)$) {？};
\end{tikzpicture}
\end{minipage}\hfill
\begin{minipage}[c]{0.51\linewidth}
三角形の1つの辺をのばすと、のばした線と、となりの辺のあいだに角ができます（緑の角 $\angle\mathrm{ACD}$）。これを、頂点Cの\textbf{外角}（がいかく）といいます。

図で、$\angle\mathrm{B}=50°$、$\angle\mathrm{A}=70°$ のとき、緑の外角は何度でしょう。
\end{minipage}

\par\vspace{4pt}
内角の和は180°なので、$\angle\mathrm{ACB}=180-(50+70)=60$ で60°。$\angle\mathrm{ACB}$ と外角は合わせて一直線（180°）なので、外角は $180-60=120$ で120°です。

120は、$50+70$ と同じです。これはたまたまではありません。外角も、$\angle\mathrm{A}+\angle\mathrm{B}$ も、どちらも「180°から $\angle\mathrm{ACB}$ をひいた角」だからです。
\bkivbfline{\fbox{三角形の外角＝となりにない2つの内角の和}}

\bkivbmidashi{多角形の内角の和}
四角形・五角形……のように、線分で囲まれた図形を\textbf{多角形}（たかくけい）といいます。多角形の内角の和は、何度でしょう。1つの頂点から対角線を引いて、三角形に分けてみます。

\begin{bkfigure}
\begin{tikzpicture}[x=1mm, y=1mm]
  % 四角形
  \fill[bkBlue!14] (0,0) -- (24,0) -- (28,20) -- cycle;
  \fill[bkOrange!16] (0,0) -- (28,20) -- (4,22) -- cycle;
  \draw[bkInk, line width=0.9pt, line join=round] (0,0) -- (24,0) -- (28,20) -- (4,22) -- cycle;
  \draw[bkInk, line width=0.6pt] (0,0) -- (28,20);
  \node[anchor=north, font=\sffamily\footnotesize, inner sep=0pt] at (14,-3) {四角形：三角形2つ};
  % 五角形
  \begin{scope}[shift={(42,0)}]
    \fill[bkBlue!14] (0,0) -- (24,0) -- (31,15) -- cycle;
    \fill[bkOrange!16] (0,0) -- (31,15) -- (14,27) -- cycle;
    \fill[bkGreen!16] (0,0) -- (14,27) -- (-4,16) -- cycle;
    \draw[bkInk, line width=0.9pt, line join=round] (0,0) -- (24,0) -- (31,15) -- (14,27) -- (-4,16) -- cycle;
    \draw[bkInk, line width=0.6pt] (0,0) -- (31,15) (0,0) -- (14,27);
    \node[anchor=north, font=\sffamily\footnotesize, inner sep=0pt] at (13,-3) {五角形：三角形3つ};
  \end{scope}
  % 六角形
  \begin{scope}[shift={(92,0)}]
    \fill[bkBlue!14] (0,0) -- (20,0) -- (30,14) -- cycle;
    \fill[bkOrange!16] (0,0) -- (30,14) -- (22,28) -- cycle;
    \fill[bkGreen!16] (0,0) -- (22,28) -- (2,28) -- cycle;
    \fill[bkPurple!14] (0,0) -- (2,28) -- (-8,14) -- cycle;
    \draw[bkInk, line width=0.9pt, line join=round] (0,0) -- (20,0) -- (30,14) -- (22,28) -- (2,28) -- (-8,14) -- cycle;
    \draw[bkInk, line width=0.6pt] (0,0) -- (30,14) (0,0) -- (22,28) (0,0) -- (2,28);
    \node[anchor=north, font=\sffamily\footnotesize, inner sep=0pt] at (11,-3) {六角形：三角形4つ};
  \end{scope}
\end{tikzpicture}
\bkfigcaption{1つの頂点から対角線を引くと、多角形は三角形に分かれる。}
\end{bkfigure}

三角形の内角を全部たすと、もとの多角形の内角を全部たしたものになります（三角形の角が、ちょうど多角形の角を分けあっているから）。表にしてみましょう。

\begin{bkcenter}
\begin{bktable}{colspec={Q[l,m,wd=32mm] *{6}{Q[c,m,wd=11mm]}}, rowsep=1.5pt, colsep=4pt}
多角形 & 三角形 & 四角形 & 五角形 & 六角形 & 七角形 & 八角形 \\
頂点の数 $n$ & 3 & 4 & 5 & 6 & 7 & 8 \\
分かれた三角形の数 & 1 & 2 & 3 & 4 & 5 & 6 \\
\SetCell{bg=bkTheme!14} 内角の和 & 180° & 360° & 540° & 720° & 900° & 1080° \\
\end{bktable}
\bktabcaption{三角形の数は、いつも頂点の数より2少ない。}
\end{bkcenter}

三角形の数が、頂点の数 $n$ より2少ないのは、選んだ頂点と、そのとなりの2つの頂点には、対角線が引けないからです（対角線は、となり合わない頂点を結ぶ線）。$n$ 個の頂点のうち、対角線の行き先になるのは $n-3$ 個で、対角線 $n-3$ 本で $n-2$ 個の三角形に分かれます。だから、
\bkivbfline{\fbox{$n$ 角形の内角の和＝$180°\times(n-2)$}}
たとえば十角形なら、$180\times(10-2)=180\times8=1440$ で1440°です。

辺の長さがすべて等しく、角の大きさもすべて等しい多角形を、\textbf{正多角形}（せいたかくけい）といいます。正 $n$ 角形の1つの内角は、内角の和を $n$ 等分して求めます。正六角形なら $720\div6=120$ で120°です。

\bkivbmidashi{多角形の外角の和は360°}
多角形の外角は、各頂点で、辺を1本ずつ同じ回り方でのばしてつくります。どの頂点でも、内角と外角を合わせると一直線の180°です。$n$ 角形では、内角と外角の組が $n$ 組あるので、
\begin{align*}
\text{外角の和}&=180°\times n-(\text{内角の和})\\
&=180°\times n-180°\times(n-2)\\
&=180°\times n-180°\times n+180°\times2=360°
\end{align*}
（2行めから3行めは、かっこをはずした（3-8）。$-180°\times(n-2)$ は、$-180°\times n$ と $+180°\times2$ になる。）$n$ が消えたので、\textbf{外角の和は、何角形でも360°}です。

\par\addvspace{4pt}\noindent\begin{minipage}[c]{0.44\linewidth}\centering
\begin{tikzpicture}[x=1mm, y=1mm]
  % 五角形のまわりを左回りに歩く。各頂点で、外角だけ左に曲がる
  \coordinate (V1) at (0,0); \coordinate (V2) at (30,0); \coordinate (V3) at (38,18);
  \coordinate (V4) at (18,32); \coordinate (V5) at (-4,18);
  \bkivbang{(V2)}{0}{66.0}{5}{bkGreen}
  \bkivbang{(V3)}{66.0}{145.0}{5}{bkGreen}
  \bkivbang{(V4)}{145.0}{212.5}{5}{bkGreen}
  \bkivbang{(V5)}{212.5}{282.5}{5}{bkGreen}
  \bkivbang{(V1)}{282.5}{360}{5}{bkGreen}
  \draw[bkInk, line width=1pt, line join=round] (V1) -- (V2) -- (V3) -- (V4) -- (V5) -- cycle;
  \draw[bkInk, line width=0.5pt, dashed] (V2) -- ++(0:10) (V3) -- ++(66:10) (V4) -- ++(145:10) (V5) -- ++(212.5:10) (V1) -- ++(282.5:10);
  \draw[-{Stealth[length=4pt]}, bkOrange, line width=1.2pt] (8,2.5) -- (20,2.5);
  \node[font=\sffamily\footnotesize, text=bkOrange!80!black, inner sep=1pt] at (17,13) {左回りに歩く};
\end{tikzpicture}
\end{minipage}\hfill
\begin{minipage}[c]{0.53\linewidth}
\textbf{歩いて確かめる}\par
五角形の辺の上を、左回りに歩きます。頂点に着くたびに、左へ向きを変えて、次の辺へ進みます。向きを変える角度が、その頂点の外角（緑）です。

1周して出発点にもどると、はじめと同じ向きを向いています。だから、とちゅうで向きを変えた角度の合計は、ちょうど1回転の360°です（1-20）。何角形でも、1周すれば1回転です。
\end{minipage}

\par\vspace{4pt}
この考えで、三角形の内角の和をもう一度確かめられます。三角形の外角の和は360°で、内角と外角の組は3組で $180\times3=540$。だから内角の和は $540-360=180$ で180°。証明と同じ答えになりました。

\bksubsection{5}{定義・公理・定理・法則}

\label{bk:4-1:teigi}④の証明では、「平行線の錯角は等しい」を使いました。それは、③で「同位角は等しい」と「対頂角は等しい」から説明しました。では、「同位角は等しい」は、何から説明できるのでしょう。

「なぜ？」「それはこうだから」「では、それはなぜ？」……と、さかのぼる問いは、いつまでも続きます。どこかで「ここから出発する」と決めないと、何も説明できません。数学では、出発点を2種類用意します。

\begin{itemize}[leftmargin=1.2\zw, itemsep=3pt, topsep=2pt, label={\color{bkTheme!80!black}●}]
\item \textbf{定義}（ていぎ）：ことばの意味を決める約束。\\
例：「二等辺三角形とは、2つの辺の長さが等しい三角形のこと」「平行とは、同じ平面の上の2本の直線が交わらないこと」\bkidx{ていぎ}{定義}
\item \textbf{公理}（こうり）：証明しないで、正しいと認めて出発点にすること。だれもが「それはそうだ」と思える、ごく当たり前のことを選びます。\\
例：「2つの点を通る直線は、1本だけある」「平行線の同位角は等しい」\bkidx{こうり}{公理}
\end{itemize}

そして、定義と公理から、証明によって正しいと示されたことを、\textbf{定理}（ていり）といいます。「対頂角は等しい」「三角形の内角の和は180°」は定理です。\bkidx{ていり}{定理}

\begin{bkfigure}
\begin{tikzpicture}[x=1mm, y=1mm,
  bl/.style={draw=#1!80!black, line width=0.9pt, fill=#1!12, rounded corners=1.5mm, minimum height=8.5mm, align=center, font=\sffamily\small, inner sep=1.5mm}]
  \node[bl=bkGray, minimum width=58mm] (d) at (30,0) {\textbf{定義}\quad 平行・角・三角形の意味};
  \node[bl=bkGray, minimum width=58mm] (k) at (94,0) {\textbf{公理}\quad 平行線の同位角は等しい};
  \node[bl=bkBlue, minimum width=64mm] (t1) at (62,14) {\textbf{定理}\quad 対頂角は等しい};
  \node[bl=bkBlue, minimum width=64mm] (t2) at (62,28) {\textbf{定理}\quad 平行線の錯角は等しい};
  \node[bl=bkOrange, minimum width=64mm] (t3) at (62,42) {\textbf{定理}\quad 三角形の内角の和は180°};
  \node[bl=bkOrange, minimum width=64mm] (t4) at (62,56) {\textbf{定理}\quad $n$ 角形の内角の和、外角の和};
  \node[anchor=west, font=\sffamily\footnotesize, text=bkGray!60!black, align=left] at (98,35) {証明で\\1段ずつ\\積み上げる};
  \draw[-{Stealth[length=4pt]}, bkGray!60!black, line width=1pt] (96,22) -- (96,50);
\end{tikzpicture}
\bkfigcaption{定義と公理を土台にして、定理を1段ずつ積み上げる。上の段は、下の段を使って証明する。}
\end{bkfigure}

建物と同じで、土台がしっかりしていれば、その上にいくらでも高く積めます。そして、いちど証明された定理は、次の証明の材料として、安心して使えます。この章の題名「きまりから考えを積み上げる」は、このことです。

\bkivbmidashi{法則}
理科や物理には、もう1つ、\textbf{法則}（ほうそく）ということばがあります。法則は、自然を実験や観察で調べて見つけたきまりです。
\begin{itemize}[leftmargin=1.2\zw, itemsep=2pt, topsep=2pt, label={\color{bkTheme!80!black}●}]
\item ばねののびは、引く力に比例する（$F=kx$、2-15）。実験で見つけた法則です。引っぱりすぎると、成り立たなくなります。
\item 物に力がはたらくと、力に比例した加速度が生まれる（力＝質量×加速度、$F=ma$）。物理の土台になる法則です（5-5、物理シリーズ）。
\end{itemize}
法則は、どれだけたくさん実験しても、「次の実験でも必ず成り立つ」と証明することはできません。じっさい、$F=ma$ は、光の速さに近い、とても速い物では少しずれることが、のちの実験でわかりました。法則は、もっと正確な実験で、なおされることがあるのです。

一方、「速さ＝距離÷時間」（2-3）は、法則ではありません。「速さということばを、こういう意味で使おう」と決めた約束、つまり定義です。だから実験でくつがえることはありません。

\begin{bkcenter}
\begin{bktable}{colspec={Q[l,m,wd=15mm] Q[l,m,wd=36mm] Q[l,m,wd=30mm] Q[l,m,wd=38mm]}, rowsep=2pt, row{1}={font=\sffamily\bfseries\small}, cells={font=\small}}
 & 何か & 正しさの決め方 & 例 \\
定義 & ことばの意味を決める約束 & 約束なので、正しい・まちがいはない & 二等辺三角形、平行、速さ＝距離÷時間 \\
公理 & 証明しないで認める出発点 & 当たり前のことを選んで、認める & 2点を通る直線は1本、平行線の同位角は等しい \\
定理 & 証明で示されたこと & 証明（定義と公理から） & 対頂角は等しい、内角の和は180° \\
法則 & 自然の中に見つけたきまり & 実験と観察（なおされることもある） & $F=kx$、$F=ma$ \\
\end{bktable}
\bktabcaption{定義・公理・定理・法則のちがい。この本では、ずっとこの使い方をする。}
\end{bkcenter}

\begin{taisetsu}
定義と公理を土台に、証明で定理を積み上げるのが数学。\par
実験と観察で法則を見つけ、なおしていくのが物理。
\end{taisetsu}

物理の法則も、いくつかの法則を出発点にして、数学の証明と同じように、たくさんのきまりを導きます。たとえば第6巻の6-15では、ばねも電気も熱も、少数の法則から同じ形の式で書けることを見ます。物理シリーズでは、$F=ma$ などの法則から、小学校から高校までの物理の式を、1つずつ導いていきます。

\bksubsection{6}{三角形が1つに決まる条件：合同条件}

\label{bk:4-1:goudou}友だちに電話で三角形を伝えて、同じ三角形をかいてもらいます。何を伝えれば、友だちは必ず、自分のと合同な三角形をかけるでしょう。

\begin{bkfigure}
\begin{tikzpicture}[x=1mm, y=1mm]
  % ① 3辺（5, 4, 3 cm を 6 mm/cm）
  \draw[bkBlue, line width=1.3pt] (0,0) -- (30,0);
  \draw[bkBlue, line width=1.3pt] (30,0) -- (19.2,14.4);
  \draw[bkBlue, line width=1.3pt] (0,0) -- (19.2,14.4);
  \node[anchor=north, font=\sffamily\footnotesize, inner sep=1pt] at (15,-0.5) {\qty{5}{cm}};
  \node[anchor=south west, font=\sffamily\footnotesize, inner sep=1pt] at (25.4,7.6) {\qty{3}{cm}};
  \node[anchor=south east, font=\sffamily\footnotesize, inner sep=1pt] at (9.2,7.6) {\qty{4}{cm}};
  \node[anchor=north, font=\sffamily\bfseries\small, inner sep=0pt] at (17,-6) {① 3辺};
  % ② 2辺とそのあいだの角
  \begin{scope}[shift={(46,0)}]
    \bkivbang{(0,0)}{0}{40}{6}{bkOrange}
    \draw[bkBlue, line width=1.3pt] (0,0) -- (30,0);
    \draw[bkBlue, line width=1.3pt] (0,0) -- (40:24);
    \draw[bkInk, line width=0.6pt, dashed] (30,0) -- (40:24);
    \node[anchor=north, font=\sffamily\footnotesize, inner sep=1pt] at (15,-0.5) {\qty{5}{cm}};
    \node[anchor=south east, font=\sffamily\footnotesize, inner sep=1pt] at (40:12) {\qty{4}{cm}};
    \node[font=\sffamily\footnotesize, inner sep=0.5pt] at (20:10) {40°};
    \node[anchor=north, font=\sffamily\bfseries\small, inner sep=0pt] at (17,-6) {② 2辺とそのあいだの角};
  \end{scope}
  % ③ 1辺とその両はしの角
  \begin{scope}[shift={(92,0)}]
    \bkivbang{(0,0)}{0}{40}{6}{bkOrange}
    \bkivbang{(30,0)}{120}{180}{5}{bkOrange}
    \draw[bkBlue, line width=1.3pt] (0,0) -- (30,0);
    \draw[bkInk, line width=0.6pt, dashed] (0,0) -- (20.21,16.96) -- (30,0);
    \node[anchor=north, font=\sffamily\footnotesize, inner sep=1pt] at (15,-0.5) {\qty{5}{cm}};
    \node[font=\sffamily\footnotesize, inner sep=0.5pt] at (20:10) {40°};
    \node[font=\sffamily\footnotesize, inner sep=0.5pt] at ($(30,0)+(150:9.5)$) {60°};
    \node[anchor=north, font=\sffamily\bfseries\small, inner sep=0pt] at (17,-6) {③ 1辺とその両はしの角};
  \end{scope}
\end{tikzpicture}
\bkfigcaption{青い線と色のついた角を伝えれば、点線の部分は自然に決まる。}
\end{bkfigure}

\begin{enumerate}[label={\bkcircled[bkDeep]{\arabic*}}, leftmargin=2.4\zw, itemsep=2pt, topsep=2pt]
\item \textbf{3辺}：「\qty{5}{cm}、\qty{4}{cm}、\qty{3}{cm}」。\qty{5}{cm} の辺を引き、両はしからコンパスで\qty{4}{cm} と\qty{3}{cm} の円をかくと、交わる点が3つめの頂点です。上と下に1つずつ交わりますが、どちらでかいても、裏返すと重なる（合同）三角形です。
\item \textbf{2辺とそのあいだの角}：「\qty{5}{cm} と\qty{4}{cm} で、そのあいだの角が40°」。2本の辺の向きと長さが決まるので、2つのはしを結べば三角形が1つに決まります。
\item \textbf{1辺とその両はしの角}：「\qty{5}{cm} で、両はしの角が40°と60°」。両はしから出る2本の線の向きが決まるので、2本が交わる点が1つに決まります。
\end{enumerate}

この3つを、三角形の\textbf{合同条件}といいます。2つの三角形で、①②③のどれかが等しければ、2つは合同です。合同条件は、「三角形が1つに決まるための条件」なのです。

\bkivbmidashi{これだけでは決まらない}
\par\addvspace{4pt}\noindent\begin{minipage}[c]{0.4\linewidth}\centering
\begin{tikzpicture}[x=1mm, y=1mm]
  \draw[bkBlue, line width=1pt, fill=bkBlue!10, line join=round] (0,0) -- (16,0) -- (11.4,9.6) -- cycle;
  \draw[bkOrange, line width=1pt, fill=bkOrange!12, line join=round] (22,0) -- (54,0) -- (44.8,19.2) -- cycle;
  \node[anchor=north, font=\sffamily\footnotesize, inner sep=0pt] at (27,-2.5) {3つの角は同じ。大きさがちがう};
\end{tikzpicture}
\end{minipage}\hfill
\begin{minipage}[c]{0.57\linewidth}
「3つの角が40°、60°、80°」だけでは、形は決まっても、大きさが決まりません。小さい三角形も、大きい三角形もかけます。3つの角が等しい三角形は、合同とはかぎらないのです（このような「形が同じ」三角形は、4-3の相似で学びます）。
\end{minipage}

\par\vspace{6pt}
\par\addvspace{4pt}\noindent\begin{minipage}[c]{0.46\linewidth}\centering
\begin{tikzpicture}[x=1mm, y=1mm]
  % B=40°、AB=5 cm、AC=4 cm（1 cm を 8 mm でかく）
  \fill[bkBlue!10] (0,0) -- (49.69,0) -- (30.64,25.71) -- cycle;
  \fill[bkOrange!20] (0,0) -- (11.59,0) -- (30.64,25.71) -- cycle;
  \bkivbang{(0,0)}{0}{40}{6}{bkGray}
  \draw[bkGray, line width=0.5pt, dashed] (30.64,25.71) ++(-140:32) arc (-140:-40:32);
  \draw[bkInk, line width=1pt] (0,0) -- (54,0);
  \draw[bkInk, line width=1pt] (0,0) -- (30.64,25.71);
  \draw[bkBlue, line width=1pt] (30.64,25.71) -- (49.69,0);
  \draw[bkOrange, line width=1pt] (30.64,25.71) -- (11.59,0);
  \node[anchor=north east, font=\sffamily\small, inner sep=1.5pt] at (0,0) {B};
  \node[anchor=south, font=\sffamily\small, inner sep=1.5pt] at (30.64,26) {A};
  \node[anchor=north east, font=\sffamily\small, inner sep=1pt] at (11.0,-0.4) {C};
  \node[anchor=north west, font=\sffamily\small, inner sep=1pt] at (50.3,-0.4) {C$'$};
  \node[font=\sffamily\footnotesize, inner sep=0.5pt] at (20:10) {40°};
  \node[anchor=south east, font=\sffamily\footnotesize, inner sep=1pt] at (14,13) {\qty{5}{cm}};
\end{tikzpicture}
\end{minipage}\hfill
\begin{minipage}[c]{0.51\linewidth}
「AB＝\qty{5}{cm}、AC＝\qty{4}{cm}、$\angle\mathrm{B}=40°$」も、だめです。40°は、\qty{5}{cm} と\qty{4}{cm} の辺の「あいだの角」ではありません。Aを中心に半径\qty{4}{cm} の円をかくと、Bから出た線と2か所（CとC$'$）で交わり、三角形が2つかけてしまいます。BCは約\qty{1.4}{cm}か約\qty{6.2}{cm}です。合同条件の②で「あいだの角」と言うのは、このためです。
\end{minipage}

\par\vspace{4pt}
4-0のロボットのうでを思い出しましょう。肩・ひじ・手を結ぶと、三角形ができます。上うで\qty{30}{cm}、前うで\qty{20}{cm}、肩から手までのまっすぐな距離の3辺が決まっているので、合同条件①から、三角形の形は1つに決まります。ただし、裏返した形（ひじが上とひじが下）もかけるので、答えは2通りでした。三角形の辺や角を、合同条件ごとに計算で求める方法（正弦定理・余弦定理）は、4-7で学びます。

\bksubsection{7}{逆と反例}

\label{bk:4-1:gyaku}「$\triangle\mathrm{ABC}$ が正三角形ならば、$\triangle\mathrm{ABC}$ は二等辺三角形である」は正しい文です。正三角形は3つの辺が等しいので、とくに2つの辺も等しいからです。この文の「ならば」の前の部分を\textbf{仮定}、後ろの部分を\textbf{結論}といいます（④の証明と同じことば）。

仮定と結論を入れかえた文を、もとの文の\textbf{逆}（ぎゃく）といいます。
\bkivbfline{もとの文\enspace「正三角形ならば、二等辺三角形」\\[2pt]\hspace*{1.5\zw}逆\enspace「二等辺三角形ならば、正三角形」}
逆は正しいでしょうか。辺の長さが\qty{5}{cm}、\qty{5}{cm}、\qty{3}{cm} の三角形は、2つの辺が等しいので二等辺三角形です。でも3つの辺は等しくないので、正三角形ではありません。だから逆は正しくありません。

\par\addvspace{4pt}\noindent\begin{minipage}[c]{0.44\linewidth}\centering
\begin{tikzpicture}[x=1mm, y=1mm]
  \draw[bkBlue, line width=1pt, fill=bkBlue!8] (0,0) ellipse (27 and 15);
  \draw[bkOrange, line width=1pt, fill=bkOrange!14] (-8,-1) ellipse (13 and 8.5);
  \node[font=\sffamily\footnotesize, text=bkBlue!80!black, inner sep=0pt] at (0,10.5) {二等辺三角形};
  \node[font=\sffamily\footnotesize, text=bkOrange!80!black, inner sep=0pt] at (-8,-1) {正三角形};
  \fill[bkRed] (14,-2) circle (1);
  \node[anchor=north, font=\sffamily\footnotesize, inner sep=1pt] at (14,-3) {5, 5, 3};
\end{tikzpicture}
\end{minipage}\hfill
\begin{minipage}[c]{0.53\linewidth}
正三角形は、二等辺三角形のなかま（内側の円）です。だから「正三角形ならば二等辺三角形」は正しい。でも、二等辺三角形のなかには、正三角形でないもの（赤い点）があります。

逆が正しくないことを示すには、「仮定にあてはまるのに、結論にあてはまらない例」を1つ見つければ十分です。この例を\textbf{反例}（はんれい）といいます。
\end{minipage}

\par\vspace{4pt}
逆が正しいこともあります。③の「平行ならば、同位角は等しい」の逆「同位角が等しいならば、平行」は正しい文でした。もとの文が正しくても、逆は正しいとはかぎりません。1つずつ確かめる必要があります。

数でも同じです。「$x=3$ ならば $x^2=9$」は正しい。逆の「$x^2=9$ ならば $x=3$」は正しくありません。反例は $x=-3$ です。$(-3)^2=9$ なのに、$x=3$ ではないからです（3-10：2乗して9になる数は3と $-3$）。

\begin{seikaku}
\textbf{反例は1つでよい、証明は全部の場合}\par
「正しくない」と言うには、反例が1つあれば十分です。でも「正しい」と言うには、たくさんの例で確かめるだけでは足りません（どこかに反例があるかもしれない）。全部の場合で成り立つことを、証明する必要があります。\bkhatenaref で三角形をいくつはかっても「必ず」と言えなかったのは、このためです。直角三角形についての定理と、その逆は、4-4（三平方の定理）で使います。
\end{seikaku}

\begin{kotoba}
\small
\bkword{定義・公理}{ことばの意味の約束。証明しないで認める出発点。}
\bkword{定理・証明}{定義と公理から、理由をつないで正しいと示したこと。その文章。}
\bkword{法則}{実験と観察で見つけた自然のきまり。なおされることもある。}
\bkword{合同}{ずらす・回す・裏返すで、ぴったり重なること。}
\bkword{対頂角・同位角・錯角}{対頂角はいつも等しい。平行線の同位角・錯角は等しい。}
\bkword{内角・外角}{$n$ 角形の内角の和は $180°\times(n-2)$、外角の和は360°。}
\bkword{逆・反例}{仮定と結論を入れかえた文。正しくないことを示す1つの例。}
\end{kotoba}

%% ---------------- 例題 ----------------
\bksection{reidai}{例題でたしかめよう}

\begin{reidai}[平行線と角]
\noindent\begin{minipage}[c]{0.56\linewidth}
右の図で、$\ell\parallel m$ です。$\angle a=65°$ のとき、$\angle x$、$\angle y$、$\angle z$ の大きさを求めましょう。
\end{minipage}\hfill
\begin{minipage}[c]{0.4\linewidth}\centering
\begin{tikzpicture}[x=1mm, y=1mm]
  \bkivbang{(14,20)}{0}{65}{5}{bkBlue}
  \bkivbang{(14,20)}{180}{245}{5}{bkOrange}
  \bkivbang{(4.67,0)}{0}{65}{5}{bkOrange}
  \bkivbang{(4.67,0)}{65}{180}{4}{bkGreen}
  \draw[bkInk, line width=0.9pt] (-8,20) -- (36,20) (-8,0) -- (36,0);
  \draw[bkInk, line width=0.9pt] (0.47,-9) -- (18.2,29);
  \node[anchor=west, font=\sffamily\small, inner sep=1pt] at (36.5,20) {$\ell$};
  \node[anchor=west, font=\sffamily\small, inner sep=1pt] at (36.5,0) {$m$};
  \node[font=\sffamily\small, inner sep=0.5pt] at ($(14,20)+(32:8.5)$) {$a$};
  \node[font=\sffamily\small, inner sep=0.5pt] at ($(14,20)+(212:8.5)$) {$x$};
  \node[font=\sffamily\small, inner sep=0.5pt] at ($(4.67,0)+(32:8.5)$) {$y$};
  \node[font=\sffamily\small, inner sep=0.5pt] at ($(4.67,0)+(125:7.5)$) {$z$};
\end{tikzpicture}
\end{minipage}\par
\kangaekata 対頂角は等しい。平行線の同位角・錯角は等しい。一直線は180°。\par
\shiki $\angle x$ は $\angle a$ の対頂角なので\quad $\angle x=65°$\par
\hspace*{2.6\zw}$\angle y$ は $\angle a$ の同位角で、$\ell\parallel m$ なので\quad $\angle y=65°$\par
\hspace*{2.6\zw}$\angle z$ と $\angle y$ を合わせると一直線なので\quad $\angle z=180-65=115$\par
\kotae $\angle x=65°$、$\angle y=65°$、$\angle z=115°$
\end{reidai}
$\angle x$ と $\angle y$ は錯角です。錯角は等しいので、ここからも $\angle y=\angle x=65°$ と確かめられます。

\begin{jibun}
例題1の図で、$\angle a=48°$ のとき、$\angle x$、$\angle y$、$\angle z$ の大きさを求めましょう。
\bkjans{$\angle x=48°$（対頂角）、$\angle y=48°$（同位角）、$\angle z=180-48=132$ で132°}
\end{jibun}

\begin{reidai}[三角形の内角と外角]
$\triangle\mathrm{ABC}$ で、$\angle\mathrm{A}=50°$、$\angle\mathrm{B}=70°$ です。$\angle\mathrm{C}$ と、頂点Cの外角の大きさを求めましょう。\par
\kangaekata 内角の和は180°。外角は、となりにない2つの内角の和（④）。\par
\shiki $\angle\mathrm{C}=180-(50+70)=180-120=60$\par
\hspace*{2.6\zw}外角$=50+70=120$\quad（確かめ：$180-60=120$）\par
\kotae $\angle\mathrm{C}=60°$、外角120°
\end{reidai}

\begin{jibun}
$\triangle\mathrm{ABC}$ で、$\angle\mathrm{A}=35°$、$\angle\mathrm{B}=80°$ です。$\angle\mathrm{C}$ と、頂点Cの外角の大きさを求めましょう。
\bkjans{$\angle\mathrm{C}=180-(35+80)=65$ で65°。外角は $35+80=115$ で115°}
\end{jibun}

\begin{reidai}[二等辺三角形の角]
$\mathrm{AB}=\mathrm{AC}$ の二等辺三角形ABCで、頂角 $\angle\mathrm{A}=40°$ です。底角 $\angle\mathrm{B}$ と $\angle\mathrm{C}$ の大きさを求めましょう。\par
\kangaekata 二等辺三角形は線対称なので、2つの底角は等しい（②）。内角の和は180°。\par
\shiki 2つの底角の和は\quad $180-40=140$\par
\hspace*{2.6\zw}2つは等しいので、1つは\quad $140\div2=70$\par
\kotae $\angle\mathrm{B}=\angle\mathrm{C}=70°$
\end{reidai}

\begin{jibun}
$\mathrm{AB}=\mathrm{AC}$ の二等辺三角形ABCで、頂角 $\angle\mathrm{A}=110°$ です。底角の大きさを求めましょう。
\bkjans{$(180-110)\div2=70\div2=35$ で35°}
\end{jibun}

\begin{reidai}[多角形の内角と外角]
（1）正八角形の内角の和と、1つの内角の大きさを求めましょう。\par
（2）五角形の4つの外角が70°、80°、65°、75°です。残りの外角は何度ですか。\par
（3）1つの内角が150°の正多角形は、正何角形ですか。\par
\kangaekata 内角の和は $180°\times(n-2)$。外角の和は何角形でも360°。内角と外角を合わせると180°。\par
\shiki（1）$180\times(8-2)=180\times6=1080$。1つの内角は $1080\div8=135$\par
\hspace*{2.6\zw}（2）$360-(70+80+65+75)=360-290=70$\par
\hspace*{2.6\zw}（3）1つの外角は $180-150=30$。外角の和は360°なので\par
\hspace*{4.6\zw}$n=360\div30=12$\par
\kotae（1）和1080°、1つの内角135°\quad（2）70°\quad（3）正十二角形
\end{reidai}
（3）は、内角の和の式からも確かめられます。正十二角形の内角の和は $180\times10=1800$ で、$1800\div12=150$。合っています。

\begin{jibun}
（1）正十角形の内角の和と、1つの内角の大きさを求めましょう。\par
（2）六角形の5つの外角が50°、70°、60°、55°、65°です。残りの外角は何度ですか。\par
（3）1つの内角が140°の正多角形は、正何角形ですか。
\bkjans{（1）\nobreak $180\times8=1440$ で1440°、$1440\div10=144$ で144°。（2）\nobreak $360-300=60$ で60°。（3）\nobreak 外角 $180-140=40$、$360\div40=9$ で正九角形}
\end{jibun}

\begin{reidai}[平行線のあいだの折れ線]
下の左の図で、$\ell\parallel m$ です。$\angle x$ の大きさを求めましょう。
\begin{bkcenter}
\begin{tikzpicture}[x=1mm, y=1mm]
  \foreach \sx/\aux in {0/0, 58/1} {
    \begin{scope}[shift={(\sx,0)}]
      {\ifnum\aux=1
        \bkivbang{(21.45,12)}{140}{180}{5}{bkBlue}
        \bkivbang{(21.45,12)}{180}{215}{6}{bkOrange}
        \draw[bkGreen!60!black, line width=0.9pt, dash pattern=on 3pt off 1.5pt] (-6,12) -- (34,12);
        \node[anchor=west, font=\sffamily\footnotesize, text=bkGreen!50!black, inner sep=1pt] at (34.5,12) {引いた平行線};
      \else
        \bkivbang{(21.45,12)}{140}{215}{4.5}{bkGreen}
        \node[font=\sffamily\small, inner sep=0.5pt] at ($(21.45,12)+(177.5:7.5)$) {$x$};
      \fi
      \bkivbang{(0,30)}{-40}{0}{6}{bkBlue}
      \bkivbang{(4.31,0)}{0}{35}{6}{bkOrange}
      \draw[bkInk, line width=0.9pt] (-6,30) -- (34,30) (-6,0) -- (34,0);
      \draw[bkInk, line width=0.9pt] (0,30) -- (21.45,12) -- (4.31,0);
      \node[anchor=west, font=\sffamily\small, inner sep=1pt] at (34.5,30) {$\ell$};
      \node[anchor=west, font=\sffamily\small, inner sep=1pt] at (34.5,0) {$m$};
      \node[font=\sffamily\footnotesize, inner sep=0.5pt] at ($(0,30)+(-20:10.5)$) {40°};
      \node[font=\sffamily\footnotesize, inner sep=0.5pt] at ($(4.31,0)+(17.5:10.5)$) {35°};
      \node[anchor=west, font=\sffamily\small, inner sep=1pt] at (22.5,13.5) {P};}
    \end{scope}}
  \node[font=\sffamily\footnotesize, text=bkBlue!80!black, inner sep=0.5pt] at ($(79.45,12)+(160:9)$) {40°};
  \node[font=\sffamily\footnotesize, text=bkOrange!80!black, inner sep=0.5pt] at ($(79.45,12)+(197.5:10)$) {35°};
\end{tikzpicture}
\end{bkcenter}
\kangaekata 上の右の図のように、Pを通って $\ell$ に平行な直線を引くと、$\angle x$ が上と下の2つに分かれる。上の部分は $\ell$ との錯角、下の部分は $m$ との錯角。\par
\shiki 上の部分＝40°（$\ell$ と、引いた線の錯角）\par
\hspace*{2.6\zw}下の部分＝35°（$m$ と、引いた線の錯角。引いた線は $m$ にも平行）\par
\hspace*{2.6\zw}$\angle x=40+35=75$\par
\kotae 75°
\end{reidai}
平行線のあいだの角は、「平行線を1本ふやす」と、知っている錯角に分けられます。答えは、両はしの角をたしたものになります。

\begin{jibun}
例題5の図で、40°が32°、35°が46°のとき、$\angle x$ の大きさを求めましょう。
\bkjans{Pを通って $\ell$ に平行な線を引く。錯角で $32+46=78$ で78°}
\end{jibun}

\begin{reidai}[合同条件]
次のア〜エで、$\triangle\mathrm{ABC}$ と $\triangle\mathrm{DEF}$ が必ず合同になるものを選び、そのときの合同条件を答えましょう。\par
ア\enspace $\mathrm{AB}=\mathrm{DE}=\qty{5}{cm}$、$\mathrm{BC}=\mathrm{EF}=\qty{7}{cm}$、$\angle\mathrm{B}=\angle\mathrm{E}=40°$\par
イ\enspace $\angle\mathrm{A}=\angle\mathrm{D}=60°$、$\angle\mathrm{B}=\angle\mathrm{E}=50°$、$\angle\mathrm{C}=\angle\mathrm{F}=70°$\par
ウ\enspace $\mathrm{BC}=\mathrm{EF}=\qty{6}{cm}$、$\angle\mathrm{B}=\angle\mathrm{E}=45°$、$\angle\mathrm{C}=\angle\mathrm{F}=75°$\par
エ\enspace $\mathrm{AB}=\mathrm{DE}=\qty{6}{cm}$、$\mathrm{AC}=\mathrm{DF}=\qty{5}{cm}$、$\angle\mathrm{B}=\angle\mathrm{E}=50°$\par
\kangaekata 等しい辺と角が、合同条件①②③のどれかにあてはまるかを見る。角は、どの辺の「あいだ」や「両はし」かを確かめる。\par
\shiki ア：$\angle\mathrm{B}$ は、辺ABと辺BCのあいだの角。②にあてはまる\par
\hspace*{2.6\zw}イ：角だけで辺の長さがない。大きさが決まらない\par
\hspace*{2.6\zw}ウ：$\angle\mathrm{B}$ と $\angle\mathrm{C}$ は、辺BCの両はしの角。③にあてはまる\par
\hspace*{2.6\zw}エ：ABとACのあいだの角は $\angle\mathrm{A}$。$\angle\mathrm{B}$ は「あいだの角」ではない。\par
\hspace*{4.6\zw}⑥の図のように、三角形が2つかけることがある\par
\kotae アは「2辺とそのあいだの角」、ウは「1辺とその両はしの角」で合同。イとエは合同とはかぎらない
\end{reidai}

\begin{jibun}
次のア〜ウで、2つの三角形が必ず合同になるものを選び、合同条件を答えましょう。\par
ア\enspace $\mathrm{AB}=\mathrm{DE}=\qty{4}{cm}$、$\mathrm{BC}=\mathrm{EF}=\qty{6}{cm}$、$\mathrm{CA}=\mathrm{FD}=\qty{5}{cm}$\par
イ\enspace $\mathrm{AB}=\mathrm{DE}=\qty{5}{cm}$、$\mathrm{AC}=\mathrm{DF}=\qty{4}{cm}$、$\angle\mathrm{B}=\angle\mathrm{E}=40°$\par
ウ\enspace $\mathrm{AC}=\mathrm{DF}=\qty{3}{cm}$、$\angle\mathrm{A}=\angle\mathrm{D}=50°$、$\angle\mathrm{C}=\angle\mathrm{F}=60°$
\bkjans{アは「3辺」、ウは「1辺とその両はしの角」で合同。イは、$\angle\mathrm{B}$ が2辺のあいだの角ではないので、合同とはかぎらない（三角形が2つかける）}
\end{jibun}

\begin{reidai}[逆と反例]
次の文の逆を書き、逆が正しいかどうかを答えましょう。正しくないときは、反例を1つ答えましょう。\par
（1）$x=2$ ならば $x^2=4$\par
（2）四角形が正方形ならば、4つの辺の長さが等しい。\par
\kangaekata 仮定と結論を入れかえる。逆の仮定にあてはまるのに、結論にあてはまらない例をさがす。\par
\shiki（1）逆は「$x^2=4$ ならば $x=2$」。$x=-2$ でも $(-2)^2=4$ だが、$x=2$ ではない\par
\hspace*{2.6\zw}（2）逆は「4つの辺の長さが等しい四角形ならば、正方形」。\par
\hspace*{4.6\zw}ひし形（4つの辺が等しく、角が直角でない四角形）は、正方形ではない\par
\kotae（1）正しくない。反例 $x=-2$\quad（2）正しくない。反例は、角が直角でないひし形
\end{reidai}

\begin{jibun}
次の文の逆を書き、正しいかどうかを答えましょう。正しくないときは、反例を1つ答えましょう。\par
（1）$x=5$ ならば $x^2=25$\par
（2）ある整数が6の倍数ならば、その整数は2の倍数である。
\bkjans{（1）\nobreak 逆「$x^2=25$ ならば $x=5$」。正しくない。反例 $x=-5$。（2）\nobreak 逆「2の倍数ならば6の倍数」。正しくない。反例は4（2の倍数だが6の倍数でない）}
\end{jibun}

%% ---------------- まちがえやすいところ ----------------
\bksection{tsumazuki}{まちがえやすいところをチェック}

\begin{tsumazuki}{$\angle\mathrm{ABC}$ の頂点}
\bkng{$\angle\mathrm{ABC}$ は、頂点Aの角}
\bkok{$\angle\mathrm{ABC}$ は、まん中に書いたBが頂点の角}
3つの文字のまん中が頂点です。$\angle\mathrm{BAC}$ ならAの角、$\angle\mathrm{ACB}$ ならCの角です。
\end{tsumazuki}

\begin{tsumazuki}{平行かどうかを確かめずに、同位角を等しいとする}
\bkng{2本の直線に1本の直線が交わったので、同位角は等しい}
\bkok{2本の直線が平行なときだけ、同位角は等しい}
③の図の左のように、平行でない2本では、同位角の大きさはちがいます。問題文や図の「$\ell\parallel m$」を確かめてから使いましょう。
\end{tsumazuki}

\begin{tsumazuki}{$n$ 角形の内角の和を $180°\times n$ にする}
\bkng{五角形の内角の和は $180\times5=900$ で900°}
\bkok{五角形の内角の和は $180\times(5-2)=540$ で540°}
三角形に分けると、三角形の数は頂点の数より2少なくなります。三角形（$n=3$）で $180\times(3-2)=180$ になるか、確かめるくせをつけましょう。
\end{tsumazuki}

\begin{tsumazuki}{逆も正しいと思いこむ}
\bkng{「正方形ならば4辺が等しい」は正しいので、「4辺が等しいならば正方形」も正しい}
\bkok{逆が正しいかどうかは、別に確かめる。ひし形が反例}
もとの文が正しくても、逆が正しいとはかぎりません。反例をさがしてみましょう。
\end{tsumazuki}

\begin{tsumazuki}{たくさん確かめたので「必ず」と言う}
\bkng{三角形を10こはかって、どれも180°に近かったから、必ず180°}
\bkok{必ず成り立つと言うには、証明が必要}
はかった例がいくつあっても、ほかの三角形で成り立つとはかぎりません。どんな場合にも使える理由（証明）をさがしましょう。
\end{tsumazuki}

\begin{tsumazuki}{「2辺と1つの角」なら合同、と思いこむ}
\bkng{2辺と1つの角が等しいので、合同}
\bkok{2辺と、その2辺の「あいだの角」が等しいときだけ、合同}
角が2辺のあいだにないと、三角形が2つかけることがあります（⑥の図）。角が、どの2辺にはさまれているかを、図に書きこんで確かめましょう。
\end{tsumazuki}

%% ---------------- 物理・工学ではこう使う ----------------
\begin{butsuri}
\bkivbmidashi{三角形は、形が変わらない：トラス}
\par\addvspace{4pt}\noindent\begin{minipage}[c]{0.55\linewidth}
4本の棒を角で回るようにつないだ四角形の枠は、横から押すとつぶれます。4辺の長さだけでは、形が1つに決まらないからです。三角形の枠は、3辺で形が1つに決まる（合同条件①）ので、つぶれません。
\end{minipage}\hfill
\begin{minipage}[c]{0.42\linewidth}\centering
\begin{tikzpicture}[x=1mm, y=1mm]
  % 四角形の枠と、つぶれた形
  \draw[bkGray, line width=1.5pt, line join=round] (0,0) -- (14,0) -- (14,14) -- (0,14) -- cycle;
  \draw[bkOrange, line width=1.2pt, dashed, line join=round] (0,0) -- (14,0) -- (21,12.1) -- (7,12.1) -- cycle;
  \foreach \x/\y in {0/0,14/0,14/14,0/14} {\filldraw[fill=white, draw=bkGray!80!black, line width=0.5pt] (\x,\y) circle (1);}
  \draw[-{Stealth[length=4pt]}, bkRed, line width=1pt] (-6.5,14) -- (-1.8,14);
  % 三角形の枠
  \begin{scope}[shift={(28,0)}]
    \draw[bkBlue, line width=1.5pt, line join=round] (0,0) -- (16,0) -- (8,13.9) -- cycle;
    \foreach \x/\y in {0/0,16/0,8/13.9} {\filldraw[fill=white, draw=bkBlue!80!black, line width=0.5pt] (\x,\y) circle (1);}
    \draw[-{Stealth[length=4pt]}, bkRed, line width=1pt] (1.5,13.9) -- (6.2,13.9);
  \end{scope}
  \node[anchor=north, font=\sffamily\footnotesize, inner sep=0pt] at (10,-2.5) {つぶれる};
  \node[anchor=north, font=\sffamily\footnotesize, inner sep=0pt] at (36,-2.5) {つぶれない};
\end{tikzpicture}
\end{minipage}
\par\vspace{4pt}
鉄橋や鉄塔は、細い鋼材を三角形に組んで作ります（\textbf{トラス}）。家の壁の中の筋かい（柱のあいだのななめの木）も、三角形を作って、地震で家がゆがむのを防ぎます。

\end{butsuri}

\begin{butsuri}
\bkivbmidashi{ロボット掃除機は、1周で360°回る}
\par\addvspace{4pt}\noindent\begin{minipage}[c]{0.6\linewidth}
ロボット掃除機には、向きの変わり方をはかるセンサー（ジャイロセンサー。半導体の小さな部品）が入っていて、曲がった角度をたし合わせて、いまの向きを知ります。

L字形の部屋のかべにそって1周すると、左に90°曲がる角が5か所、右に90°曲がる角が1か所あります。右に曲がるのを反対の向きのマイナス（2-2）とすると、
\end{minipage}\hfill
\begin{minipage}[c]{0.37\linewidth}\centering
\begin{tikzpicture}[x=1mm, y=1mm]
  \fill[bkGray!12] (0,0) -- (40,0) -- (40,22) -- (22,22) -- (22,38) -- (0,38) -- cycle;
  \draw[bkGray!80!black, line width=1.6pt, line join=round] (0,0) -- (40,0) -- (40,22) -- (22,22) -- (22,38) -- (0,38) -- cycle;
  \draw[-{Stealth[length=4pt]}, bkOrange, line width=1pt, rounded corners=2.5mm] (6,4) -- (36,4) -- (36,18) -- (18,18) -- (18,34) -- (4,34) -- (4,8);
  \foreach \x/\y in {32.6/7.4, 32.6/14.6, 14.6/30.6, 7.4/30.6, 9.2/9.8} {\node[font=\sffamily\footnotesize, text=bkBlue!80!black, inner sep=0pt] at (\x,\y) {左};}
  \node[font=\sffamily\footnotesize, text=bkRed!80!black, inner sep=0pt] at (21.4,14.6) {右};
  \filldraw[fill=bkNavy!70, draw=bkNavy, line width=0.5pt] (8,4) circle (2.4);
\end{tikzpicture}
\end{minipage}
\[ 90\times5+(-90)=450-90=360 \]
で、ちょうど1回転です。出発点にもどったとき合計が360°からずれていれば、センサーのずれがわかり、直すことができます。
\end{butsuri}

\begin{butsuri}
\bkivbmidashi{単位の「定義」}
\qty{1}{m}の長さは、「光が、何もない空間（真空）を $\dfrac{1}{\num{299792458}}$ 秒のあいだに進む距離」と決められています。これは実験で見つけた法則ではなく、世界じゅうで同じ長さを使うための約束、つまり定義です。光の速さが「秒速\qty{299792458}{m}」ぴったりなのは、はかった結果ではなく、そう定義したからです。
\end{butsuri}

\begin{column}{三角形の内角の和が、180°にならない世界}
今から約2300年前、ギリシャのユークリッドは、それまでに知られていた図形の性質を、少数の定義と公理から証明で積み上げた本『原論』をまとめました。その公理の1つが、平行線についての公理です。

19世紀になって、数学者たちは、この平行線の公理を別のものに取りかえても、矛盾（むじゅん：つじつまが合わないこと）のない図形の世界が作れることに気づきました。たとえば、地球のような球の表面です。北極から経度0°の線にそって赤道まで進み、赤道にそって経度90°の所まで進み、北極へもどると、3つの角がどれも90°の三角形ができます。内角の和は270°です。球の表面では、平らな面の公理が成り立たないからです。

公理を変えると、定理も変わります。この本の図形は、平らな面の公理の上に積み上げたものです。地球の大きさにくらべて小さい範囲（学校の校庭など）なら、平らな面と考えて問題ありません。
\end{column}

%% ---------------- 章末問題 ----------------
\bkendproblems
\bkproblemgroup{基本}{1}

\bkproblem 例題1の図で、$\angle a=72°$ のとき、$\angle x$、$\angle y$、$\angle z$ の大きさを求めましょう（\bkreidairef{1}）。
\bkpans{$\angle x=72°$、$\angle y=72°$、$\angle z=180-72=108$ で108°}

\bkproblem $\triangle\mathrm{ABC}$ で、$\angle\mathrm{A}=62°$、$\angle\mathrm{B}=43°$ です。$\angle\mathrm{C}$ と、頂点Cの外角の大きさを求めましょう（\bkreidairef{2}）。
\bkpans{$\angle\mathrm{C}=180-105=75$ で75°。外角は $62+43=105$ で105°}

\bkproblem $\mathrm{AB}=\mathrm{AC}$ の二等辺三角形ABCで、頂角 $\angle\mathrm{A}=36°$ です。底角の大きさを求めましょう（\bkreidairef{3}）。
\bkpans{$(180-36)\div2=144\div2=72$ で72°}

\bkproblem 次の問いに答えましょう（\bkreidairef{4}）。\\
（1）正五角形の内角の和と、1つの内角の大きさを求めましょう。\\
（2）七角形の6つの外角が、どれも50°です。残りの外角は何度ですか。\\
（3）1つの内角が160°の正多角形は、正何角形ですか。
\bkpans{（1）\nobreak $180\times3=540$ で540°、$540\div5=108$ で108°。（2）\nobreak $360-50\times6=360-300=60$ で60°。（3）\nobreak 外角20°、$360\div20=18$ で正十八角形}

\bkproblem 例題5の図の40°が38°、35°が44°のときの、$\angle x$ を求めましょう（\bkreidairef{5}）。
\bkpans{$38+44=82$ で82°}

\bkproblem 次のア・イで、2つの三角形が必ず合同になるものを選び、合同条件を答えましょう（\bkreidairef{6}）。\\
ア\enspace $\mathrm{AB}=\mathrm{DE}=\qty{8}{cm}$、$\angle\mathrm{A}=\angle\mathrm{D}=30°$、$\mathrm{AC}=\mathrm{DF}=\qty{6}{cm}$\\
イ\enspace $\angle\mathrm{A}=\angle\mathrm{D}=90°$、$\angle\mathrm{B}=\angle\mathrm{E}=30°$、$\angle\mathrm{C}=\angle\mathrm{F}=60°$
\bkpans{アは「2辺とそのあいだの角」で合同（$\angle\mathrm{A}$ はABとACのあいだの角）。イは角だけなので、合同とはかぎらない}

\bkproblemgroup{チャレンジ}{2}

\bkproblem 次の文の逆を書き、正しいかどうかを答えましょう。正しくないときは、反例を1つ答えましょう（\bkreidairef{7}）。\\
（1）$x=4$ ならば $x^2=16$\qquad（2）整数が10の倍数ならば、5の倍数である。
\bkpans{（1）\nobreak 逆「$x^2=16$ ならば $x=4$」。正しくない。反例 $x=-4$。（2）\nobreak 逆「5の倍数ならば10の倍数」。正しくない。反例は5（または15など）}

\bkproblem 正 $n$ 角形の1つの外角が24°です。$n$ と、内角の和を求めましょう（\bkreidairef{4}）。
\bkpans{$n=360\div24=15$。内角の和は $180\times(15-2)=180\times13=2340$ で2340°}

\bkproblem ロボット掃除機が五角形の部屋のかべにそって左回りに1周し、4つの角で左に70°、75°、80°、65°曲がりました。残りの角（？）では、左に何度曲がりましたか（\bkreidairef{4}、④の「歩いて確かめる」）。
\begin{bkcenter}
\begin{tikzpicture}[x=1mm, y=1mm]
  \coordinate (V0) at (0,0); \coordinate (V1) at (40,0); \coordinate (V2) at (47.52,20.67);
  \coordinate (V3) at (22.95,37.88); \coordinate (V4) at (-3.98,10.95);
  \fill[bkGray!12] (V0) -- (V1) -- (V2) -- (V3) -- (V4) -- cycle;
  \foreach \v/\a/\b in {V1/0/70, V2/70/145, V3/145/225, V4/225/290, V0/290/360} {
    {\bkivbang{(\v)}{\a}{\b}{6}{bkGreen}
     \draw[bkInk, line width=0.5pt, dashed] (\v) -- ++(\a:11);}}
  \draw[bkGray!80!black, line width=1.4pt, line join=round] (V0) -- (V1) -- (V2) -- (V3) -- (V4) -- cycle;
  \node[font=\sffamily\footnotesize, inner sep=0.5pt] at ($(V1)+(35:10)$) {70°};
  \node[font=\sffamily\footnotesize, inner sep=0.5pt] at ($(V2)+(107.5:10)$) {75°};
  \node[font=\sffamily\footnotesize, inner sep=0.5pt] at ($(V3)+(185:10)$) {80°};
  \node[font=\sffamily\footnotesize, inner sep=0.5pt] at ($(V4)+(257.5:10)$) {65°};
  \node[font=\sffamily\bfseries\footnotesize, text=bkRed!80!black, inner sep=0.5pt] at ($(V0)+(325:10)$) {？};
  \draw[-{Stealth[length=4pt]}, bkOrange, line width=1pt] (8,4) -- (24,4);
  \filldraw[fill=bkNavy!70, draw=bkNavy, line width=0.5pt] (6,4) circle (2);
  \node[font=\sffamily\footnotesize, text=bkOrange!80!black, inner sep=1pt] at (22,17) {左回りに1周};
\end{tikzpicture}
\bkfigcaption{緑の角が、それぞれの角で左に曲がる角（外角）}
\end{bkcenter}
\bkpans{1周で360°回るので、$360-(70+75+80+65)=360-290=70$ で70°}

\begin{furikaeri}
\begin{checklist}
\item 平行線の同位角・錯角から、角を求められる。\bkfurref{bk:4-1:heikou}
\item 内角の和が180°のわけと、多角形の角の和がわかる。\bkfurref{bk:4-1:naikaku}
\item 定義・公理・定理・法則のちがいを説明でき、合同条件・逆・反例を使える。\bkfurref{bk:4-1:teigi}
\end{checklist}
\end{furikaeri}

\bkfukidashi{sora}{はかるだけじゃ「必ず」とは言えないけど、証明すれば、まだはかっていない三角形のことまでわかるんだね。}
